% ENEM 2020-2025 — 11 questões MCQ — Análise Combinatória I
% Fonte: lista "Análise Combinatória I" do site projetoagathaedu.com.br
% Omitidas: Q5 (alternativas com figuras), Q7 (imagem essencial), Q9 (imagem essencial), Q14 (alternativas com figuras)
% Nota: Q8 tinha alternativas corrompidas no HTML; reconstruídas pelo MathJax e validadas pelo gabarito

---
tipo: Múltipla Escolha
dificuldade: Média
nivel: médio
disciplina: Matemática
assunto: Análise Combinatória
gabarito: A
tags: [senha, aplicativo, caracteres, probabilidade]
fonte: concurso
concurso: ENEM
banca: INEP
ano: 2025
numero: 1
---
\question
Ao realizar o cadastro em um aplicativo de investimentos, foi solicitado ao usuário que criasse uma senha, sendo permitido o uso somente dos seguintes caracteres:

-- algarismos de 0 a 9;

-- 26 letras minúsculas do alfabeto;

-- 26 letras maiúsculas do alfabeto;

-- 6 caracteres especiais: !, @, \#, \$, *, \&.

Três tipos de estruturas para senha foram apresentadas ao usuário:

-- tipo I: formada por quaisquer quatro caracteres distintos, escolhidos dentre os permitidos;

-- tipo II: formada por cinco caracteres distintos, iniciando por três letras, seguidas por um algarismo e, ao final, um caractere especial;

-- tipo III: formada por seis caracteres distintos, iniciando por duas letras, seguidas por dois algarismos e, ao final, dois caracteres especiais.

Considerando \( P_1 \), \( P_2 \) e \( P_3 \), as probabilidades de se descobrirem ao acaso, na primeira tentativa, as senhas dos tipos I, II e III, respectivamente.

Nessas condições, o tipo de senha que apresenta a menor probabilidade de ser descoberta ao acaso, na primeira tentativa, é
\begin{choices}
\correctchoice tipo I, pois \( P_1 < P_2 < P_3 \).
\choice tipo I, pois tem menor quantidade de caracteres.
\choice tipo II, pois tem maior quantidade de letras.
\choice tipo III, pois \( P_3 < P_2 < P_1 \).
\choice tipo III, pois tem maior quantidade de caracteres.
\end{choices}

---
tipo: Múltipla Escolha
dificuldade: Fácil
nivel: médio
disciplina: Matemática
assunto: Análise Combinatória
gabarito: B
tags: [campeonato, turno e returno, times, equação]
fonte: concurso
concurso: ENEM-PPL
banca: INEP
ano: 2021
numero: 2
---
\question
Um diretor esportivo organiza um campeonato no qual haverá disputa de times em turno e returno, isto é, cada time jogará duas vezes com todos os outros, totalizando 380 partidas a serem disputadas.

A quantidade de times \( x \) que faz parte desse campeonato pode ser calculada pela equação
\begin{oneparchoices}
\choice \( x = 380 - x^2 \)
\correctchoice \( x^2 - x = 380 \)
\choice \( x^2 = 380 \)
\choice \( 2x - x = 380 \)
\choice \( 2x = 380 \)
\end{oneparchoices}

---
tipo: Múltipla Escolha
dificuldade: Média
nivel: médio
disciplina: Matemática
assunto: Análise Combinatória
gabarito: B
tags: [futebol, grupos, calendário, descanso]
fonte: concurso
concurso: ENEM
banca: INEP
ano: 2025
numero: 3
---
\question
Foram convidadas 32 equipes para um torneio de futebol, que foram divididas em 8 grupos com 4 equipes, sendo que, dentro de um grupo, cada equipe disputa uma única partida contra cada uma das demais equipes de seu grupo. A primeira e a segunda colocadas de cada grupo seguem para realizar as 8 partidas da próxima fase do torneio, chamada oitavas de final. Os vencedores das partidas das oitavas de final seguem para jogar as 4 partidas das quartas de final. Os vencedores das quartas de final disputam as 2 partidas das semifinais, e os vencedores avançam para a grande final, que define a campeã do torneio.

Pelas regras do torneio, cada equipe deve ter um período de descanso de, no mínimo, 3 dias entre dois jogos por ela disputados, ou seja, se um time disputar uma partida, por exemplo, num domingo, só poderá disputar a partida seguinte a partir da quinta-feira da mesma semana.

O número mínimo de dias necessários para a realização desse torneio é
\begin{oneparchoices}
\choice 22
\correctchoice 25
\choice 28
\choice 48
\choice 64
\end{oneparchoices}

---
tipo: Múltipla Escolha
dificuldade: Fácil
nivel: médio
disciplina: Matemática
assunto: Análise Combinatória
gabarito: B
tags: [montadora, carros, configurações, opcionais]
fonte: concurso
concurso: ENEM
banca: INEP
ano: 2025
numero: 4
---
\question
Uma montadora de automóveis divulgou que oferta a seus clientes mais de 1.000 configurações diferentes de carro, variando o modelo, a motorização, os opcionais e a cor do veículo. Atualmente, ela oferece 7 modelos de carros com 2 tipos de motores: 1.0 e 1.6. Já em relação aos opcionais, existem 3 escolhas possíveis: central multimídia, rodas de liga leve e bancos de couro, podendo o cliente optar por incluir um, dois, três ou nenhum dos opcionais disponíveis.

Para ser fiel à divulgação feita, a quantidade mínima de cores que a montadora deverá disponibilizar a seus clientes é
\begin{oneparchoices}
\choice 8
\correctchoice 9
\choice 11
\choice 18
\choice 24
\end{oneparchoices}

---
tipo: Múltipla Escolha
dificuldade: Fácil
nivel: médio
disciplina: Matemática
assunto: Análise Combinatória
gabarito: A
tags: [hotel, quartos, numeração, algarismo]
fonte: concurso
concurso: ENEM
banca: INEP
ano: 2025
numero: 6
---
\question
Um hotel de 3 andares está sendo construído. Cada andar terá 100 quartos. Os quartos serão numerados de 100 a 399 e cada um terá seu número afixado à porta. Cada número será composto por peças individuais, cada uma simbolizando um único algarismo.

Qual a quantidade mínima de peças, simbolizando o algarismo 2, necessárias para identificar o número de todos os quartos?
\begin{oneparchoices}
\correctchoice 160
\choice 157
\choice 130
\choice 120
\choice 60
\end{oneparchoices}

---
tipo: Múltipla Escolha
dificuldade: Média
nivel: médio
disciplina: Matemática
assunto: Permutações Simples
gabarito: E
tags: [anagrama, Harry Potter, vogais, consoantes]
fonte: concurso
concurso: ENEM
banca: INEP
ano: 2025
numero: 8
---
\question
Nos livros \textit{Harry Potter}, um anagrama do nome do personagem “TOM MARVOLO RIDDLE” gerou a frase “I AM LORD VOLDEMORT”.

Suponha que Harry quisesse formar todos os anagramas da frase “I AM POTTER”, de tal forma que as vogais e consoantes aparecessem sempre intercaladas, e sem considerar o espaçamento entre as letras.

Nessas condições, o número de anagramas formados é dado por
\begin{oneparchoices}
\choice \( 9! \)
\choice \( 4! \cdot 5! \)
\choice \( 2 \cdot 4! \cdot 5! \)
\choice \( \frac{9!}{2} \)
\correctchoice \( \frac{4! \cdot 5!}{2} \)
\end{oneparchoices}

---
tipo: Múltipla Escolha
dificuldade: Média
nivel: médio
disciplina: Matemática
assunto: Permutações Simples
gabarito: D
tags: [e-mail, anagrama, Eduardo, bloco]
fonte: concurso
concurso: ENEM-Digital
banca: INEP
ano: 2020
numero: 10
---
\question
Eduardo deseja criar um e-mail utilizando um anagrama exclusivamente com as sete letras que compõem o seu nome, antes do símbolo @. O e-mail terá a forma *******@site.com.br e será de tal modo que as três letras “edu” apareçam sempre juntas e exatamente nessa ordem. Ele sabe que o e-mail eduardo@site.com.br já foi criado por outro usuário e que qualquer outro agrupamento das letras do seu nome forma um e-mail que ainda não foi cadastrado. De quantas maneiras Eduardo pode criar um e-mail desejado?
\begin{oneparchoices}
\choice 59
\choice 60
\choice 118
\correctchoice 119
\choice 120
\end{oneparchoices}

---
tipo: Múltipla Escolha
dificuldade: Fácil
nivel: médio
disciplina: Matemática
assunto: Divisibilidade
gabarito: C
tags: [brincadeira, múltiplos, sequência, futebol]
fonte: concurso
concurso: ENEM-Digital
banca: INEP
ano: 2020
numero: 11
---
\question
“1, 2, 3, GOL, 5, 6, 7, GOL, 9, 10, 11, GOL, 13, GOL, 15, GOL, 17, 18, 19, GOL, 21, 22, 23, GOL, 25, ...”

Para a Copa do Mundo de Futebol de 2014, um bar onde se reuniam amigos para assistir aos jogos criou uma brincadeira. Um dos presentes era escolhido e tinha que dizer, numa sequência em ordem crescente, os números naturais não nulos, trocando os múltiplos de 4 e os números terminados em 4 pela palavra GOL. A brincadeira acabava quando o participante errava um termo da sequência.

Um dos participantes conseguiu falar até o número 103, respeitando as regras da brincadeira.

O total de vezes em que esse participante disse a palavra GOL foi
\begin{oneparchoices}
\choice 20
\choice 28
\correctchoice 30
\choice 35
\choice 40
\end{oneparchoices}

---
tipo: Múltipla Escolha
dificuldade: Fácil
nivel: médio
disciplina: Matemática
assunto: Análise Combinatória
gabarito: B
tags: [celular, senha, toques, códigos]
fonte: concurso
concurso: ENEM-Digital
banca: INEP
ano: 2020
numero: 12
---
\question
Um modelo de telefone celular oferece a opção de desbloquear a tela usando um padrão de toques como senha.

Os toques podem ser feitos livremente nas 4 regiões numeradas da tela, sendo que o usuário pode escolher entre 3, 4 ou 5 toques ao todo.

Qual expressão representa o número total de códigos existentes?
\begin{oneparchoices}
\choice \( 4^5 - 4^4 - 4^3 \)
\correctchoice \( 4^5 + 4^4 + 4^3 \)
\choice \( 4^5 \cdot 4^4 \cdot 4^3 \)
\choice \( (4!)^5 \)
\choice \( 4^5 \)
\end{oneparchoices}

---
tipo: Múltipla Escolha
dificuldade: Fácil
nivel: médio
disciplina: Matemática
assunto: Combinações
gabarito: E
tags: [flores, canteiro, trio, variedades]
fonte: concurso
concurso: ENEM-PPL
banca: INEP
ano: 2020
numero: 13
---
\question
A prefeitura de uma cidade está renovando os canteiros de flores de suas praças. Entre as possíveis variedades que poderiam ser plantadas, foram escolhidas cinco: amor-perfeito, cravina, petúnia, margarida e lírio. Em cada um dos canteiros, todos com composições diferentes, serão utilizadas somente três variedades distintas, não importando como elas serão dispostas.

Um funcionário deve determinar os trios de variedades de flores que irão compor cada canteiro.

De acordo com o disposto, a quantidade de trios possíveis é dada por
\begin{oneparchoices}
\choice 5
\choice \( 5 \cdot 3 \)
\choice \( \frac{5!}{(5-3)!} \)
\choice \( \frac{5!}{(5-3)!2!} \)
\correctchoice \( \frac{5!}{(5-3)!3!} \)
\end{oneparchoices}

---
tipo: Múltipla Escolha
dificuldade: Fácil
nivel: médio
disciplina: Matemática
assunto: Análise Combinatória
gabarito: C
tags: [campeonato, pontos corridos, empates, futebol]
fonte: concurso
concurso: ENEM-PPL
banca: INEP
ano: 2020
numero: 15
---
\question
Um determinado campeonato de futebol, composto por 20 times, é disputado no sistema de pontos corridos. Nesse sistema, cada time joga contra todos os demais times em dois turnos, isto é, cada time joga duas partidas com cada um dos outros times, sendo que cada jogo pode terminar empatado ou haver um vencedor.

Sabendo-se que, nesse campeonato, ocorreram 126 empates, o número de jogos em que houve ganhador é igual a
\begin{oneparchoices}
\choice 64
\choice 74
\correctchoice 254
\choice 274
\choice 634
\end{oneparchoices}
